%%%%%%%%%%%%%%%%%%%%%%%%%%%%%%%%%%%%%%%%%%%%%%%%%%%%%%%%%%%%%%%%%%%%%
%% Optical Thermometry in Monolayer WS2: Raman and PL Spectroscopy, 5 K to 1273 K
%% Draft v1 — G. Lee
%%%%%%%%%%%%%%%%%%%%%%%%%%%%%%%%%%%%%%%%%%%%%%%%%%%%%%%%%%%%%%%%%%%%%
\documentclass[letterpaper]{article}

\usepackage[T1]{fontenc}
\usepackage[utf8]{inputenc}

\usepackage{geometry}
\geometry{margin=1in}
\usepackage{setspace}
% \doublespacing

\usepackage[style=chem-acs]{biblatex}
\addbibresource{references.bib}

\usepackage{graphicx}
\usepackage{float}
\usepackage{amsmath}
\usepackage{amssymb}

\usepackage{chemformula}
\usepackage[version=4]{mhchem}

\setcounter{secnumdepth}{-1}

\usepackage{authblk}

% TODO: confirm author order with co-authors
\author[1]{G. Lee}
\author[1]{A. Borel}
\author[2]{T. Taniguchi}
\author[3]{K. Watanabe}
\author[1]{F. Sirotti}
\author[1]{F. Cadiz}

\affil[1]{Laboratoire de Physique de la Matière Condensée, CNRS,
Ecole polytechnique, Institut Polytechnique de Paris, 91120 Palaiseau, France}
\affil[2]{International Center for Materials Nanoarchitectonics, National Institute for
Materials Science, 1-1 Namiki, Tsukuba 305-0044, Japan}
\affil[3]{Research Center for Functional Materials, National Institute for Materials
Science, 1-1 Namiki, Tsukuba 305-0044, Japan}

\title{Optical Thermometry in Monolayer \ce{WS2}: Raman and
Photoluminescence Spectroscopy from 5~K to 1273~K}

\date{*Email: fabian.cadiz@polytechnique.edu}

\begin{document}

\maketitle

%%%%%%%%%%%%%%%%%%%%%%%%%%%%%%%%%%%%%%%%%%%%%%%%%%%%%%%%%%%%%%%%%%%%%
\begin{abstract}
%%%%%%%%%%%%%%%%%%%%%%%%%%%%%%%%%%%%%%%%%%%%%%%%%%%%%%%%%%%%%%%%%%%%%

Reliable temperature measurement is essential for
characterizing and engineering two-dimensional (2D) semiconductor devices, yet
optical thermometry of 2D materials has so far been demonstrated only over
narrow temperature windows spanning a few hundred kelvin. Here, we establish a comprehensive optical
thermometry framework for monolayer \ce{WS2} from 5~K to 1273~K.
Using two independent optical probes on
complementary samples, we show that: (i) the out-of-plane
A$_{1g}$ Raman mode shifts linearly with temperature above 100~K, with
non-encapsulated \ce{WS2} exhibiting a slope of $-0.0138~\mathrm{cm^{-1}/K}$
(fitted over 298--1223~K) and hBN-encapsulated samples a measurably smaller
slope of $-0.0121~\mathrm{cm^{-1}/K}$ (fitted over 100--1273~K);
(ii) photoluminescence spectra measured over the same 5--1273~K range provide
an independent optical indicator: the neutral-exciton emission $X_A^0$
redshifts continuously with temperature and is well described by the
semi-empirical Pässler model. The two indicators probe distinct physical
mechanisms, lattice anharmonicity and bandgap renormalization, and together
provide contact-free thermometry for monolayer \ce{WS2} from 5~K to 1273~K,
an approach applicable to 2D semiconductors more broadly.

\end{abstract}

%%%%%%%%%%%%%%%%%%%%%%%%%%%%%%%%%%%%%%%%%%%%%%%%%%%%%%%%%%%%%%%%%%%%%
\section{Introduction}
%%%%%%%%%%%%%%%%%%%%%%%%%%%%%%%%%%%%%%%%%%%%%%%%%%%%%%%%%%%%%%%%%%%%%

The isolation of two-dimensional (2D) atomic crystals, beginning with graphene
and extending to a broad family of layered compounds \cite{novoselov-2005}, has
established a class of materials whose properties are set by their reduced
dimensionality. Among them, monolayer transition metal dichalcogenides (TMDs)
are direct-gap semiconductors whose optical response is dominated by tightly
bound excitons even at room temperature \cite{mak-2010,splendiani-2010}. The
strong Coulomb interaction and reduced dielectric screening in these
atomically-thin materials give exciton binding energies of several hundred meV
and make the optical response acutely sensitive to the dielectric surroundings
\cite{ugeda-2014,chernikov-2014}.

High-temperature annealing is emerging as a route to defect engineering in 2D
semiconductors: in monolayer \ce{MoS2} it generates chalcogen-vacancy defects
with characteristic optical signatures \cite{mitterreiter-2021}, and in
hBN-encapsulated monolayer \ce{WS2} annealing near 1100~K produces ultra-narrow
emission lines that act as single-photon emitters \cite{lee-2026d}. Controlling
such processes requires reliable knowledge of the temperature of the 2D layer
itself, which cannot be obtained by attaching a contact sensor to a layer one
atom thick. The optical response of the monolayer offers a way out: both its
excitonic emission and its Raman spectrum shift measurably with temperature, so
the spectrum itself can serve as a built-in, contact-free thermometer.

Optical thermometry based on photoluminescence (PL) and Raman spectroscopy is
well established as a non-contact probe at the nanoscale \cite{brites-2012},
and has been extended to 2D materials such as graphene, where the G and 2D
Raman modes provide local temperature readout \cite{calizo-2007}. In the case of
monolayer \ce{MoS2} and \ce{WS2}, temperature-dependent Raman spectroscopy has
been used to extract thermal conductivities \cite{yan-2014,sahoo-2013}, with
linear peak shifts reported over limited temperature windows, typically
spanning a few hundred kelvin below $\sim$550~K \cite{huang-2016,lanzillo-2013}. However, the combination of (i) an ultra-wide
temperature range from cryogenic to near-dissociation temperatures, (ii)
simultaneous access to both Raman and PL observables, and (iii) a quantitative
assessment of the robustness of these indicators across different encapsulation
conditions has not been demonstrated to date.

A key practical challenge is that high-power laser illumination, required for
sufficient Raman signal in monolayer TMDs, induces strong photodoping in
non-encapsulated samples, as established for the related molybdenum
dichalcogenides \cite{cadiz-2016}, making PL inaccessible after Raman
acquisition. Encapsulation of the TMD monolayer between hexagonal boron nitride
(hBN) layers suppresses this photodoping effect \cite{cadiz-2016} and
dramatically narrows excitonic linewidths \cite{cadiz-2017,ajayi-2017}, enabling PL
spectroscopy under the same conditions as Raman. The trade-off is a reduction in
the Raman signal due to the additional hBN layers, leading to missing data points
at certain temperatures. The two measurement configurations are therefore
complementary rather than interchangeable.

In this work, we exploit this complementarity to establish an optical thermometry
framework for monolayer \ce{WS2} spanning 5~K to 1273~K---a range that, to our
knowledge, has not been covered by a single 2D material system in a single
study. Both samples are mounted on a suspended SiC micro-heater membrane inside
a cryostat, so that the monolayer can be heated locally while the surrounding
environment stays cold. The refractory character of SiC allows the membrane to
be operated repeatedly up to 1273~K, so that the same sample can be measured
over the full range without changing the optical setup. The non-encapsulated sample gives the Raman thermal response from
room temperature to 1223~K, while the hBN-encapsulated sample provides PL over
the full range together with complementary Raman data. Calibrating both
signatures against a known temperature turns the spectrum itself into the
thermometer: a single Raman or PL acquisition then reports the temperature of
the monolayer at that spot.

%%%%%%%%%%%%%%%%%%%%%%%%%%%%%%%%%%%%%%%%%%%%%%%%%%%%%%%%%%%%%%%%%%%%%
\section{Experimental Methods}
%%%%%%%%%%%%%%%%%%%%%%%%%%%%%%%%%%%%%%%%%%%%%%%%%%%%%%%%%%%%%%%%%%%%%

\subsection{Sample Preparation}

For the non-encapsulated sample (Experiment 1), monolayer \ce{WS2} was
mechanically exfoliated from a bulk crystal and transferred directly onto the
suspended SiC micro-membrane of a commercial micro-heater chip (EDEN
Instruments). For the encapsulated sample (Experiment 2), hBN/\ce{WS2}/hBN
heterostructures were assembled by dry-stamp transfer \cite{castellanos-gomez-2014}
and subsequently placed onto a second membrane chip. The SiC membrane is
120~nm thick, suspended between doped Si contacts, and enables uniform Joule
heating of the central region up to $\sim$1200~K when a dc current is applied.
The temperature of the membrane was calibrated against the applied electrical
current by the manufacturer \cite{lee-2026d}.

\subsection{Temperature Control}

Two complementary temperature stages were used for the encapsulated sample.
From 5~K to room temperature, the sample was cooled in a He-flow cryostat and
heated uniformly using the cryostat's built-in heater. In this regime, the
cryostat thermometer measures the temperature of the entire cold finger, which
is in equilibrium with the sample. From room temperature to $\sim$1273~K, the
SiC membrane was heated resistively under vacuum ($\sim$$10^{-3}$~mbar) while
the cryostat environment remained at low temperature. The current-to-temperature
conversion was obtained from the manufacturer's calibration. For the
non-encapsulated sample, only high-temperature measurements were performed (room
temperature to 1273~K under vacuum).

\subsection{Optical Measurements}

All measurements used a 514.5~nm continuous-wave laser (Cobolt). For Raman
spectroscopy, the excitation power was $\sim$1~mW; for PL spectroscopy of the
encapsulated sample, the power was reduced to $\lesssim$10~$\mu$W to avoid
photo-induced heating. The laser was focused to a $\sim$1~$\mu$m spot through a
50$\times$ long-working-distance objective. The scattered or emitted light was
collected by the same objective, filtered with volume Bragg notch filters centered
at 514.5~nm (for Raman), and dispersed by a grating spectrometer onto a
Peltier-cooled CCD. For each temperature step, the current was increased
gradually, held at the target value for 30~min, and then reduced slowly back to
zero before the next acquisition. PL spectra were fitted with a pseudo-Voigt
profile (see Supporting Information) to extract the peak energy $E_0$ and the
total full-width at half-maximum (FWHM) at each temperature.

%%%%%%%%%%%%%%%%%%%%%%%%%%%%%%%%%%%%%%%%%%%%%%%%%%%%%%%%%%%%%%%%%%%%%
\section{Results}
%%%%%%%%%%%%%%%%%%%%%%%%%%%%%%%%%%%%%%%%%%%%%%%%%%%%%%%%%%%%%%%%%%%%%

\subsection{Raman Spectroscopy as a Thermal Indicator}

\subsubsection{Non-encapsulated \ce{WS2} (RT to 1223~K)}

A representative Raman spectrum of the non-encapsulated \ce{WS2} monolayer at
room temperature is shown in Fig.~\ref{fig:fig1}(c). In addition to the
principal zone-center modes E$^1_{2g}$($\Gamma$) ($\sim$352~cm$^{-1}$) and
A$_{1g}$($\Gamma$) ($\sim$414~cm$^{-1}$), several zone-boundary combination
modes are resolved, including 2LA(M) ($\sim$349~cm$^{-1}$), 3LA(M)
($\sim$517~cm$^{-1}$), and A$_{1g}$($\Gamma$)+LA(M) ($\sim$576~cm$^{-1}$),
consistent with prior reports on monolayer \ce{WS2} \cite{berkdemir-2013}.
The A$_{1g}$($\Gamma$) mode is used as the primary thermal indicator in what
follows; the overlapping E$^1_{2g}$($\Gamma$)+2LA(M) doublet is analysed
separately in the Supporting Information (Sec.~S1), where it serves as a
secondary, qualitative check.

Figure~\ref{fig:fig2}(a) shows a two-dimensional intensity map of the Raman
spectra as a function of temperature from room temperature to 1223~K.
With increasing temperature, both peaks redshift continuously and broaden,
while their intensity decreases. The redshift arises from two contributions:
thermal expansion of the lattice, which lowers the phonon frequency through the
Grüneisen relation, and anharmonic phonon--phonon coupling, whose three- and
four-phonon decay channels become increasingly populated \cite{huang-2016}. The signal vanishes above
$\sim$1223~K, consistent with thermal dissociation of the
\ce{WS2} lattice through sulfur sublimation \cite{brainard-1968} (see also
Supporting Information).

The A$_{1g}$ peak position extracted from these spectra is shown in
Fig.~\ref{fig:fig2}(c) (red squares). A linear fit restricted to $T \geq 100$~K
yields:

\begin{equation}
\omega_{A_{1g}}(T) = \omega_0 + \left(\frac{d\omega}{dT}\right) T,
\end{equation}

with a slope $d\omega/dT = -0.0138 \pm 0.0002~\mathrm{cm^{-1}/K}$ and intercept
$\omega_0 = 417.95~\mathrm{cm^{-1}}$ ($R^2 = 0.995$, $T = 298$--$1223$~K); the
same dataset analysed with the anti-Stokes A$_{1g}$ line included gives
$-0.0134~\mathrm{cm^{-1}/K}$ \cite{lee-2026d}. A linear fit to the
A$_{1g}$ shift of monolayer \ce{WS2} on quartz by Huang \textit{et al.} gave
$-0.0122~\mathrm{cm^{-1}/K}$ over 84--543~K, although those authors restricted
their quantitative analysis of this mode to below 380~K, above which strain
from wrinkling of the monolayer on the substrate distorted the response
\cite{huang-2016}. Reported first-order coefficients for the A$_{1g}$ mode
span $-0.006$ to $-0.0149~\mathrm{cm^{-1}/K}$ across studies \cite{huang-2016},
so that values obtained on different samples, substrates and fitting ranges are
not directly comparable; the two configurations measured here were acquired on
the same platform with identical excitation and fitting procedures, making the
difference between them, quantified below, a controlled comparison. The A$_{1g}$ mode is preferred as the
primary thermometric indicator because, unlike the overlapping
E$^1_{2g}$($\Gamma$)+2LA(M) doublet, it is spectrally isolated, so that its
position can be determined unambiguously at every temperature.

\begin{figure}[H]
  \centering
  \includegraphics[width=\linewidth]{figures/fig1.png}
  \caption{\textbf{Experimental platform and reference spectra.}
    (a) Schematic of the in-situ measurement platform: hBN/\ce{WS2}/hBN
    heterostructure transferred onto the suspended SiC micro-membrane of the
    micro-heater chip, mounted inside a cryostat.
    (b) Optical image of the encapsulated \ce{WS2} flake on the SiC membrane
    under white-light illumination.
    (c) Room-temperature Raman spectrum of non-encapsulated monolayer \ce{WS2}
    ($\lambda_\mathrm{exc} = 514.5$~nm, $P = 1$~mW), including Lorentzian peak
    fits for the E$^1_{2g}$($\Gamma$), 2LA(M), A$_{1g}$($\Gamma$), and
    higher-order combination modes.
    High-power illumination induces photodoping in the absence of hBN encapsulation,
    precluding PL measurements on this sample.
    (d) PL spectrum of hBN-encapsulated \ce{WS2} at 5~K, showing the neutral
    exciton $X_A^0$, trion, and dark-exciton-related features
    ($\lambda_\mathrm{ex} = 514.5$~nm, $P = 10$~$\mu$W).}
  \label{fig:fig1}
\end{figure}

\subsubsection{hBN-encapsulated \ce{WS2} (5~K to 1273~K)}

The hBN-encapsulated sample allows simultaneous Raman and PL measurements under
identical laser conditions, since hBN encapsulation suppresses the photodoping
that precludes PL in the non-encapsulated sample \cite{cadiz-2016}.
The Raman and PL data presented in this and the following section were acquired
on the same encapsulated sample.

Figure~\ref{fig:fig2}(b) shows the two-dimensional Raman intensity map of the
encapsulated sample over the full temperature range from 5~K to
1273~K. The overall signal level is reduced compared to the non-encap
sample due to partial attenuation by the hBN layers. The map reveals a
temperature-dependent intensity redistribution between the two modes: at low
temperature (below $\sim$100~K), the E$^1_{2g}$+2LA mode dominates; at high
temperature, the A$_{1g}$ mode becomes the stronger feature; at intermediate
temperatures both modes are weak, leading to gaps in the extracted peak
positions. One factor that plausibly contributes to the redistribution is a
change in resonance condition: at 514.5~nm, monolayer \ce{WS2} sits in a
second-order Raman resonance that makes the 2LA(M) mode roughly twice as
intense as the first-order A$_{1g}$($\Gamma$) mode, an enhancement specific to
this excitation wavelength and governed by the electronic band structure
\cite{berkdemir-2013}. Because the excitation energy is fixed while the
electronic transitions of the sample redshift by $\sim$0.4~eV over the range
studied (Fig.~\ref{fig:fig3}(b)), the detuning between the two changes
continuously with temperature, and the second-order mode is expected to be far
more sensitive to this than the first-order one. This alone does not account
for the observation that \emph{both} modes weaken at intermediate temperatures,
and the physical origin of the redistribution is under investigation and will
be addressed in a forthcoming study.

Despite the missing data points at intermediate temperatures,
Fig.~\ref{fig:fig2}(c) (blue circles) shows that the A$_{1g}$ peak positions
obtained from the encapsulated sample across the full accessible range (5~K to
RT from cryogenic measurements; RT to 1273~K from micro-heater measurements)
follow the same qualitative linear trend established for the non-encapsulated
sample, indicating that the underlying phonon-softening mechanism is intrinsic
to the \ce{WS2} lattice regardless of encapsulation, although the two samples
differ measurably in thermal sensitivity, as quantified below. A linear fit
restricted to $T \geq 100$~K yields a slope
$d\omega/dT = -0.0121 \pm 0.0001~\mathrm{cm^{-1}/K}$ and intercept
$\omega_0 = 419.55~\mathrm{cm^{-1}}$ ($R^2 = 0.997$), a magnitude
$\sim$12\% smaller than in the non-encapsulated sample. Because this fit pools
cryostat-calibrated points below room temperature with micro-heater-calibrated
points above it, whereas the non-encapsulated fit uses micro-heater
calibration alone, we checked that the difference is not an artefact of
combining the two temperature scales: restricting the encapsulated fit to
$T \geq 298$~K, so that it too rests on the micro-heater calibration only,
gives $-0.0121~\mathrm{cm^{-1}/K}$, unchanged within the fit uncertainty.
This difference is in
the direction expected if the hBN layers mechanically constrain the monolayer:
the in-plane thermal expansion coefficient of hBN is negative
\cite{androulidakis-2020}, whereas that of \ce{WS2} is positive
\cite{huang-2016}, so an encapsulated monolayer cannot expand as freely as a
bare one and the thermal-expansion contribution to the phonon softening is
correspondingly reduced. Testing this quantitatively would require an
independent measurement of the residual strain and is left for future work.

\begin{figure}[H]
  \centering
  \includegraphics[width=\linewidth]{figures/fig2.png}
  \caption{\textbf{Raman thermal indicator: non-encapsulated and hBN-encapsulated
    \ce{WS2}.}
    (a) Two-dimensional map of Raman intensity versus temperature (298--1223~K)
    and Raman shift for the non-encapsulated sample. Both the E$^1_{2g}$+2LA and
    A$_{1g}$ modes redshift continuously with increasing temperature.
    (b) Two-dimensional map of Raman intensity versus temperature (5--1273~K)
    and Raman shift for the hBN-encapsulated sample. The white dashed line marks
    the boundary between the cryostat (5~K--RT) and micro-heater (RT--1273~K)
    acquisition stages. A temperature-dependent intensity redistribution between
    the E$^1_{2g}$+2LA (dominant at low $T$) and A$_{1g}$ (dominant at high $T$)
    modes is observed, with both peaks weak at intermediate temperatures.
    (c) A$_{1g}$($\Gamma$) peak position versus temperature extracted from both
    samples: non-encapsulated (red squares) and hBN-encapsulated (blue circles).
    Solid lines are linear fits restricted to $T \geq 100$~K, with slopes of
    $-0.0138~\mathrm{cm^{-1}/K}$ (non-encapsulated, $R^2 = 0.995$) and
    $-0.0121~\mathrm{cm^{-1}/K}$ (hBN-encapsulated, $R^2 = 0.997$). Raman and PL
    measurements on the encapsulated sample were performed simultaneously (see
    Fig.~\ref{fig:fig3}).}
  \label{fig:fig2}
\end{figure}

\subsection{Photoluminescence as a Thermal Indicator}

\textbf{Spectral evolution and peak energy.}
Figure~\ref{fig:fig3}(a) shows the evolution of the PL spectrum of the
hBN-encapsulated \ce{WS2} monolayer, plotted as a waterfall of six
representative spectra spanning 5~K to 1023~K. At 5~K the emission is dominated
by lower-lying excitonic complexes, and the neutral exciton $X_A^0$ at
$E_0 = 2.052$~eV, with a narrow Voigt linewidth of $\sim$5.3~meV, is one of the
weakest features of the spectrum (Fig.~\ref{fig:fig1}(d)). This reflects the
darkish nature of tungsten-based monolayers, in which the excitonic ground
state is optically inactive, so that photocarriers relax preferentially into
dark and charged states \cite{zinkiewicz-2020,zhang-2015}. With increasing
temperature, thermal activation redistributes population into the bright
$X_A^0$ state, whose integrated intensity rises by two orders of magnitude up
to $\sim$175~K (Supporting Information, Fig.~S5); at the same time the
lower-lying complexes broaden and merge, and above $\sim$300~K the spectrum
consists of a single $X_A^0$ line. On further heating this line redshifts
monotonically and broadens, while its emission efficiency falls through
thermally activated non-radiative recombination \cite{korn-2011} and the escape
of excitons from the radiative window \cite{robert-2016}.

The neutral-exciton peak energy $E_0(T)$ extracted from Voigt fits is shown in
Fig.~\ref{fig:fig3}(b). The redshift reflects the interplay of two mechanisms:
thermal lattice expansion, which reduces the orbital overlap and lowers the
bandgap, and electron--phonon coupling, which renormalizes the quasiparticle
energies \cite{passler-1999}. The data are excellently described by the
Pässler model \cite{passler-1997,passler-1999}:

\begin{equation}
E_0(T) = E_0(0) - \frac{\alpha \Theta}{2}
\left[ \sqrt[p]{1 + \left(\frac{2T}{\Theta}\right)^p} - 1 \right],
\label{eq:passler}
\end{equation}

where $E_0(0)$ is the zero-temperature bandgap, $\alpha$ is the high-temperature
slope, $\Theta$ is the mean phonon temperature, and $p$ is a dimensionless
dispersion exponent. The best-fit parameters obtained here are:
$E_0(0) = 2.0491 \pm 0.0036$~eV, $\alpha = (3.741 \pm 0.069)\times10^{-4}~\mathrm{eV/K}$,
$\Theta = 280.1 \pm 28.8$~K, $p = 4.04 \pm 2.23$ (the large relative uncertainty
on $p$ reflects its known correlation with $\Theta$ in this model)
($R^2 = 0.9964$, RMSE $= 7.9$~meV;
see Table~S1). These are consistent with previously reported values
for encapsulated \ce{WS2} \cite{lee-2026d}, with small deviations attributed to
differences in local strain and dielectric environment between samples. A
comparison of the Pässler model with the Varshni and O'Donnell--Chen models is
provided in the Supporting Information; both the Pässler and O'Donnell--Chen
models reproduce the data equivalently well, while the Varshni model
underestimates the curvature at low temperature.

The inversion of Eq.~\eqref{eq:passler} allows the temperature to be read
directly from the measured $E_0$, providing a model-based, contact-free
thermometer calibrated against the cryostat (5~K to RT) and the manufacturer's
chip calibration (RT to 1273~K). The practical lower limit of this PL
thermometer is set by the spectral resolution of the detection system. Towards
higher temperature, the $X_A^0$ peak broadens substantially (from
$\sim$5.3~meV at 5~K to $\sim$140~meV at 1273~K) and becomes progressively less
well defined, so that its position is determined with decreasing precision;
the measurements reported here extend to 1273~K.
The root-mean-square deviation of the Pässler fit residuals is
$7.9$~meV, corresponding to a temperature uncertainty of $\sim 21$~K in the
high-temperature (linear) regime, estimated as RMSE$/\alpha$; the uncertainty
is larger at cryogenic temperature, where $dE_0/dT \to 0$.

\textbf{Spectral linewidth.}
The $X_A^0$ emission broadens monotonically with temperature, the total Voigt
FWHM increasing from $\sim$5.3~meV at 5~K to $\sim$140~meV at 1273~K
(Supporting Information, Fig.~S5). The narrow low-temperature linewidth
reflects the quality of the hBN-encapsulated heterostructure
\cite{cadiz-2017,ajayi-2017}, while the progressive broadening at higher
temperature is consistent with an increasing phonon-scattering-limited
dephasing rate \cite{selig-2016,moody-2015}. The linewidth is not used as a
thermometric indicator here: it is strongly sample-dependent, and at the
highest temperatures the peak extends beyond the low-energy edge of the
fitting window, so that the width is no longer reliably determined.

\begin{figure}[H]
  \centering
  \includegraphics[width=\linewidth]{figures/fig3.png}
  \caption{\textbf{Photoluminescence thermal indicator in hBN-encapsulated monolayer
    \ce{WS2}.}
    (a) Waterfall plot of representative PL spectra (5~K to 1023~K) as a
    function of photon energy, normalized to the peak intensity at each
    temperature, overlaid on the full temperature-dependent intensity map. The
    $X_A^0$ neutral-exciton peak, visible at low temperature, redshifts and
    broadens continuously with increasing temperature.
    (b) Neutral-exciton peak energy $E_0$ extracted from Voigt fits (black
    squares, with error bars) and the Pässler model fit (blue line),
    Eq.~\eqref{eq:passler}, with best-fit parameters and fit quality annotated.}
  \label{fig:fig3}
\end{figure}

%%%%%%%%%%%%%%%%%%%%%%%%%%%%%%%%%%%%%%%%%%%%%%%%%%%%%%%%%%%%%%%%%%%%%
\section{Discussion}
%%%%%%%%%%%%%%%%%%%%%%%%%%%%%%%%%%%%%%%%%%%%%%%%%%%%%%%%%%%%%%%%%%%%%

Our results establish monolayer \ce{WS2} as a broadband optical thermometer
whose two independent indicators---Raman peak position and PL peak energy---are
applicable over different, overlapping temperature windows that together span
5~K to 1273~K.

The Raman indicator is the simpler of the two: a single linear calibration
curve suffices above 100~K. The hBN-encapsulated sample, however, shows a
measurably different calibration from the bare monolayer, its slope being
12\% smaller ($-0.0121$ vs.\ $-0.0138~\mathrm{cm^{-1}/K}$). Because the slope
is a differential quantity, this difference is independent of the absolute
wavenumber calibration, and it sets a floor on the accuracy of any transferred
calibration: using one sample's slope to read the other's temperature leaves a
systematic error of 12\%, or $\sim$120~K at 1000~K. A Raman calibration
obtained on a bare monolayer therefore cannot be carried over to an
encapsulated heterostructure without re-calibration. What this work establishes
is the range over which either calibration remains linear, from 100~K to
1273~K, substantially broader than the window previously reported for
\ce{WS2} Raman thermometry \cite{huang-2016}.

The PL thermometer relies on the neutral-exciton peak energy, for which the
Pässler model provides a physically grounded, analytic expression that can be
inverted to yield temperature from a single spectral acquisition.
Importantly, the PL thermometer extends the range to cryogenic temperatures (5~K),
which the Raman thermometer cannot easily access due to the difficulty of
distinguishing small thermal shifts in the nearly temperature-independent phonon
spectrum at low $T$.

The two indicators are complementary in a specific sense. Below $\sim$100~K the
Raman shift is too weakly temperature-dependent to resolve, whereas the PL peak
energy still tracks temperature with sub-kelvin precision; PL is therefore the
only usable indicator in the cryogenic regime. At high temperature the two are
comparable: converting the peak-position uncertainty of each fitted spectrum
into a temperature through the corresponding sensitivity ($\sigma_{E_0}/\alpha$
for PL, $\sigma_\omega/|d\omega/dT|$ for Raman) gives, averaged over all points
above 900~K, $\sim$12~K for PL and $\sim$13~K for Raman. The practical
distinction there is
not precision but sample configuration, since photodoping precludes PL on
non-encapsulated samples \cite{cadiz-2016}, leaving the Raman shift as the only
available indicator for bare monolayers.

Finally, it is useful to place this range in the context of Raman thermometry
on other 2D materials, where the published temperature windows are
considerably narrower: 100--320~K for suspended monolayer \ce{MoS2}
\cite{yan-2014}, 303--448~K for monolayer \ce{MoS2} on \ce{SiO2}
\cite{lanzillo-2013}, and 84--543~K for monolayer \ce{WS2} on quartz
\cite{huang-2016}, the last restricted to below 380~K for the A$_{1g}$ mode by
substrate-induced wrinkling. The linear behaviour reported here therefore
extends over a range several times wider than in any of these studies.
Departures from linearity are present but small: including a quadratic term
changes the fitted frequency by at most 0.4~cm$^{-1}$ out of a total excursion
of 12.5~cm$^{-1}$, so a single linear calibration is retained throughout.

The two samples reach their upper limits for different reasons. The
non-encapsulated monolayer dissociates through sulfur sublimation
\cite{brainard-1968}: its Raman signal is lost above $\sim$1223~K and the flake
is no longer visible afterwards (Supporting Information, Sec.~S6). The
encapsulated monolayer survives 1273~K intact, with both Raman and PL still
measurable, and the measurement was stopped there to protect the heater chip
rather than because the sample failed. hBN encapsulation therefore extends the
thermal stability of the monolayer as well as its spectral quality.

Two features of the platform make the extended range accessible: the suspended SiC membrane allows the monolayer to
be heated locally while the surrounding cryostat stays cold, and hBN
encapsulation sharpens the excitonic linewidth sufficiently for reliable peak
fitting down to 5~K \cite{cadiz-2017,ajayi-2017}. The large exciton binding
energy of monolayer \ce{WS2} ($\sim$300~meV) \cite{chernikov-2014} ensures a
well-defined excitonic peak at all temperatures studied, and the use of a
single excitation wavelength (514.5~nm) for both Raman and PL simplifies the
measurement.

%%%%%%%%%%%%%%%%%%%%%%%%%%%%%%%%%%%%%%%%%%%%%%%%%%%%%%%%%%%%%%%%%%%%%
\section{Conclusion}
%%%%%%%%%%%%%%%%%%%%%%%%%%%%%%%%%%%%%%%%%%%%%%%%%%%%%%%%%%%%%%%%%%%%%

We have established an optical thermometry framework for monolayer \ce{WS2}
spanning 5~K to 1273~K using two independent optical indicators. The
A$_{1g}$ Raman mode shifts linearly with temperature above 100~K, with slopes
of $-0.0138~\mathrm{cm^{-1}/K}$ for the bare monolayer (298--1223~K) and
$-0.0121~\mathrm{cm^{-1}/K}$ for the hBN-encapsulated heterostructure
(100--1273~K); the 12\% difference between them means that a calibration
obtained on one configuration cannot be transferred to the other. The
neutral-exciton PL peak energy, described by the Pässler model, provides a
second indicator over the same range, with a fit residual corresponding to
$\sim$21~K in the high-temperature regime.

Two questions remain open. The origin of the temperature-dependent intensity
redistribution between the E$^1_{2g}$+2LA doublet and the A$_{1g}$ mode, which
removes both features at intermediate temperatures in the encapsulated sample,
is still under investigation. And the upper limit reached here was set by the
heater chip rather than by the encapsulated monolayer, which survived 1273~K
intact; how much further hBN encapsulation extends the usable range remains to
be established.

%%%%%%%%%%%%%%%%%%%%%%%%%%%%%%%%%%%%%%%%%%%%%%%%%%%%%%%%%%%%%%%%%%%%%
\section*{Acknowledgements}
%%%%%%%%%%%%%%%%%%%%%%%%%%%%%%%%%%%%%%%%%%%%%%%%%%%%%%%%%%%%%%%%%%%%%

We thank QCMX for sample preparation assistance.
FC acknowledges support from the European Research Council (ERC Starting Grant
OneSPIN, Grant No.\ 101075855). We also acknowledge financial support from
DEEP2D (ANR-22-CE09-0013).

%%%%%%%%%%%%%%%%%%%%%%%%%%%%%%%%%%%%%%%%%%%%%%%%%%%%%%%%%%%%%%%%%%%%%
\section*{Data Availability}
%%%%%%%%%%%%%%%%%%%%%%%%%%%%%%%%%%%%%%%%%%%%%%%%%%%%%%%%%%%%%%%%%%%%%

The data that support the findings of this study will be made available in
Zenodo upon publication.

%%%%%%%%%%%%%%%%%%%%%%%%%%%%%%%%%%%%%%%%%%%%%%%%%%%%%%%%%%%%%%%%%%%%%
\section*{Supporting Information}
%%%%%%%%%%%%%%%%%%%%%%%%%%%%%%%%%%%%%%%%%%%%%%%%%%%%%%%%%%%%%%%%%%%%%

The Supporting Information contains: (S1) Individual versus merged-peak
fitting of the E$^1_{2g}$+2LA Raman doublet for both samples; (S2)
Temperature-dependent Raman spectra of the hBN-encapsulated sample; (S3)
Bandgap model comparison (Varshni, Pässler, O'Donnell--Chen) with fit
residuals; (S4) Voigt fitting procedure and representative fit quality; (S5)
Temperature dependence of the PL spectral linewidth and integrated $X_A^0$
intensity; (S6) \ce{WS2} monolayer after thermal
dissociation.

%%%%%%%%%%%%%%%%%%%%%%%%%%%%%%%%%%%%%%%%%%%%%%%%%%%%%%%%%%%%%%%%%%%%%
\printbibliography
%%%%%%%%%%%%%%%%%%%%%%%%%%%%%%%%%%%%%%%%%%%%%%%%%%%%%%%%%%%%%%%%%%%%%

\end{document}
